\chapter{Conclusão}

Este trabalho teve como objetivo investigar a aplicação da Architecture Analysis Design Language (AADL), em conjunto com o ecossistema do Robot Operating System 2 (ROS~2), como suporte ao projeto e à validação arquitetural de sistemas robóticos. A partir da modelagem estruturada dos componentes de \textit{hardware}, \textit{software} e dos canais de comunicação, foi possível analisar antecipadamente aspectos críticos relacionados ao desempenho, ao uso de recursos computacionais e à integração entre subsistemas, ainda nas fases iniciais do desenvolvimento.

Os resultados obtidos demonstram que a abordagem dirigida por modelos constitui uma ferramenta eficaz para a tomada de decisões arquiteturais fundamentadas. A utilização da AADL, por meio do ambiente OSATE, permitiu avaliar quantitativamente a distribuição da carga computacional entre os núcleos de processamento, o uso dos barramentos de comunicação e a disponibilidade de recursos não vinculados, evidenciando que a arquitetura proposta é coerente com as demandas do pipeline de visão computacional e não apresenta gargalos estruturais no escopo analisado.

O estudo de caso do robô Robinho evidenciou a aplicabilidade prática da abordagem proposta. Mesmo considerando uma aplicação de visão computacional de complexidade moderada, a arquitetura modelada apresentou margens adequadas de processamento, comunicação e memória, além de uma alocação eficiente dos processos mais exigentes em núcleos dedicados. Esses resultados indicam que a arquitetura é estável e possui potencial de escalabilidade para a inclusão de novos sensores, módulos de \textit{software} e requisitos temporais mais rigorosos. 

É importante ressaltar que os resultados obtidos dependem do nível de abstração adotado no modelo arquitetural e das estimativas de carga computacional utilizadas, não substituindo medições empíricas em execução real. Ainda assim, as análises realizadas são suficientes para validar arquiteturalmente o sistema no escopo proposto, cumprindo o papel de antecipar problemas, reduzir riscos de integração e orientar decisões técnicas antes da implementação física. O objetivo da análise não é prever com exatidão o comportamento temporal do sistema real, mas identificar tendências e gargalos potenciais que orientem decisões arquiteturais antes da implementação física.

Estudos correlatos apontam elevada correlação entre os resultados obtidos por meio de análises em AADL e medições realizadas em sistemas físicos, com diferenças reportadas da ordem de 0{,}02\% \cite{Senn2021}. Nesse contexto, trabalhos futuros incluem a comparação sistemática entre os valores estimados pelo OSATE e medições experimentais no robô Robinho, bem como a investigação de técnicas de geração automática de código a partir dos modelos arquiteturais, ampliando o nível de automação e a precisão das análises.

Ao demonstrar que é possível antecipar gargalos e validar decisões arquiteturais antes da implementação física, este trabalho contribui para a redução de riscos, custos e retrabalho no desenvolvimento de sistemas robóticos complexos, reforçando o papel da modelagem arquitetural como elemento central no projeto de Sistemas Ciberfísicos modernos.